\begin{statement}{E18}{Refuted}{The square theta formula for arbitrary sections}
The square-only formula in Lead~9 is false as a local period
identity for arbitrary elliptic sections and linearizations.
The exact formula for the theta expression specified there is
\eqref{eq:E-theta-general}. The assertion for the actual original
period family remains true, by its equation $\wp(\mu)=0$.
\end{statement}
\begin{proof}
Let $\Delta$ be a sufficiently small disc about $i$ in the upper
half-plane and choose the holomorphic logarithm of the nonvanishing
function $\vartheta_1(1/2,\tau)/\eta(\tau)$ which is real at $i$.
Put
\[
 \mu(\tau)=\frac12,\qquad
 v(\tau)=\frac{i}{\pi}
       \log\frac{\vartheta_1(1/2,\tau)}{\eta(\tau)}+i,
 \qquad \beta(\tau)=6v(\tau).
\]
At $\tau=i$, the product expansions give
\[
 \frac{\vartheta_1(1/2,i)}{\eta(i)}
 =2e^{-\pi/6}\prod_{n\geq1}(1+e^{-2\pi n})^2>1.
\]
Thus $\operatorname{Im}v(i)>1$, and after shrinking $\Delta$
one has $\operatorname{Im}v>0$. Define
\[
 \Lambda_\tau=\mathbb Z(1,0)+\mathbb Z(0,1)
 +\mathbb Z(\tau,1/2)+\mathbb Z(3,\beta(\tau)),
 \qquad T_\tau=\mathbb C^2/\Lambda_\tau.
\]
The real determinant of these four periods is nonzero, since
$\operatorname{Im}\tau>0$ and $\operatorname{Im}\beta>0$.
They form a locally uniformly full lattice, so the quotients form
a proper holomorphic family of compact complex two-tori over
$\Delta$. Each fibre has
$h^\bullet(\mathcal O)=(1,2,1)$ and
$h^\bullet(T)=(2,4,2)$, with zero higher groups, by its
translation-invariant Dolbeault complex.

This family also has the required elliptic line-bundle
description. Quotient first by $(1,0),(0,1),(\tau,1/2)$.
Writing $z$ for the first coordinate and $a=\exp(2\pi iz_2)$
for the second gives the complement of zero in the flat line
bundle on $\mathbb C/(\mathbb Z+\tau\mathbb Z)$ with multipliers
$1$ and $-1$. The theta quotient for the point $P=1/2$ identifies
this line bundle with $\mathcal O(P-O)$. The last lattice
generator acts by
$(z,a)\mapsto(z+3,\exp(2\pi i\beta)a)$.
Its elliptic translation is $-6P=0$ in the elliptic curve;
its scalar is a holomorphic contracting linearization. This
constructs the family from an elliptic family, a section and a
linearization explicitly.

There is a specified relation to the symmetric period matrix
$\Omega=\left(\begin{smallmatrix}\tau&\mu\\\mu&v\end{smallmatrix}\right)$.
Adjoining $(\mu,v)$ enlarges $\Lambda_\tau$ to the lattice
generated by $(1,0),(0,1),(\tau,\mu),(\mu,v)$; its index is six,
because $6(\mu,v)=(3,\beta)$ and no smaller multiple belongs
to the original real lattice. Identity on $\mathbb C^2$ induces
a finite holomorphic isogeny of degree six between these families.
Its differential is the identity on the Lie spaces and its
pullback on complex first cohomology is an isomorphism. The
Kodaira--Spencer tensors therefore have the same rank under
these explicitly given comparisons.

The identity $\vartheta_1(1/2+z,\tau)=\vartheta_2(z,\tau)$
and evenness of $\vartheta_2$ give
$\partial_\mu\log\vartheta_1(1/2,\tau)=0$. Since $\mu'=0$,
the proposed square is zero. The actual determinant, by the
proved heat-equation formula, is
\[
 \det\Omega'(i)=-\frac{\wp(1/2;i)}{4\pi^2}
 =-\frac{\vartheta_3(0,i)^4+\vartheta_4(0,i)^4}{12}\ne0.
\]
The last inequality follows from the theta products at the real
nome $e^{-\pi}$, whose factors give strictly positive theta
constants. Thus even the zero sets in the asserted formula
differ in this fully specified family. Extra holomorphic terms
in a different linearization contribute their derivatives to
the determinant; they must be retained, just as the $\wp$ term
must be retained when the section is not a zero of $\wp$.

For completeness, the direct images and total-space cohomology
in this local example are also determined. Write
$g:\mathscr T\to\Delta$ for the displayed proper family.
Translation-invariant forms, base change, and triviality of
holomorphic vector bundles on a disc give
\[
 R^qg_*\mathcal O_{\mathscr T}
   =\mathcal O_\Delta^{\binom2q},\qquad
 R^qg_*T_{\mathscr T/\Delta}
   =\mathcal O_\Delta^{2\binom2q}\quad(0\leq q\leq2),
\]
with all higher images zero. Shrink once more so that the
nonzero determinant stays nonzero. In the direct-image sequence
of $0\to T_{\mathscr T/\Delta}\to T_{\mathscr T}
\to g^*T_\Delta\to0$, the first boundary sends $1$ to the
nowhere-zero tensor $\kappa$, and the next boundary is the
rank-two isomorphism of E5. Thus
\[
 R^qg_*T_{\mathscr T}=
 \begin{cases}
  \mathcal O_\Delta^2&q=0,\\
  \mathcal O_\Delta^3&q=1,\\
  \mathcal O_\Delta&q=2,\\
  0&q>2.
 \end{cases}
\]
Since $\Delta$ is Stein, every total-space group
$H^q(\mathscr T,\mathcal F)$ for these three sheaves is exactly
the space of sections on $\Delta$ of the corresponding displayed
direct image. This specifies every cohomological degree without
assigning finite dimensions to the holomorphic-function spaces
on a disc.
\end{proof}
